%! Multiplying \& Dividing Radicals; Rationalizing — Unit 5, Lesson 5.3 — Student Packet Cover Sheet
\documentclass[10pt]{article}
\usepackage{algebra2-article}
\usepackage{algebra2-boxes}
\usepackage{ltablex}
\keepXColumns

\begin{document}

% Full-bleed forest banner
\begin{tikzpicture}[remember picture, overlay]
  \fill[forest] (current page.north west)
    rectangle ([yshift=-0.9in]current page.north east);
\end{tikzpicture}
\vspace{-0.8in}

\begin{center}
  {\color{white}\textbf{\LARGE Algebra 2: Shepherd}}\\[4pt]
  {\color{forestmid}\large\textbf{Unit 5 \quad Radical Functions}}\\[6pt]
  {\color{white}\textbf{\large Lesson 5.3 \quad Multiplying \& Dividing Radicals; Rationalizing}}
\end{center}

\vspace{0.22in}
\namedateperiod
\vspace{0.18in}

\begin{learningtargetbox}
\small
\begin{enumerate}[leftmargin=1.5em, itemsep=5pt]
  \item \ldots{} \textbf{multiply} radicals --- including distributing and FOIL --- and simplify afterward, because a product can hide a perfect power that neither factor had.
  \item \ldots{} recognize a \textbf{conjugate} and explain why multiplying a pair of conjugates always produces a \emph{rational} number.
  \item \ldots{} \textbf{rationalize} a denominator, monomial or binomial, and explain why doing so never changes the value of the expression.
\end{enumerate}
\end{learningtargetbox}

\vspace{0.14in}

\begin{tocbox}
\small
\renewcommand{\arraystretch}{1.5}
\begin{tabularx}{\linewidth}{@{} c l X r @{}}
\rowcolor{forestbg}
\textbf{\#} & \textbf{Component} & \textbf{Description} & \textbf{Score} \\
\midrule
1 & Warm-Up        & Spiral review: distributing and FOIL, $\sqrt a\cdot\sqrt a$, and multiplying by a form of $1$ & \blank{1.2cm} \\
2 & Guided Notes   & Multiplying, distributing and FOIL, conjugates, dividing, and rationalizing both kinds of denominator & \blank{1.2cm} \\
3 & Group Activity & \emph{Making It Rational} --- products, conjugate pairs, and rationalizing, by tier & \blank{1.2cm} \\
4 & Exit Ticket    & Quick check: multiply, FOIL, rationalize, explain, plus an SOL-style item & \blank{1.2cm} \\
5 & Homework       & Mixed products and quotients, rationalizing, and the golden ratio & \blank{1.2cm} \\
\midrule
  & \multicolumn{2}{r}{\textbf{Total}} & \blank{1.2cm} \\
\end{tabularx}
\end{tocbox}

\vspace{0.10in}

\begin{remindbox}
\small Yesterday you ran the product property \emph{backwards}, pulling perfect powers out. Today
you run it \textbf{forwards}: $\sqrt[n]{a}\cdot\sqrt[n]{b}=\sqrt[n]{ab}$ --- multiply the outsides
together, multiply the insides together, \emph{then} simplify. Simplifying last matters, because a
product often creates a perfect power that neither factor had ($\sqrt6\cdot\sqrt{10}=\sqrt{60}
=2\sqrt{15}$). Radicals \textbf{distribute} and \textbf{FOIL} exactly like binomials, so
$\left(\sqrt5+\sqrt2\right)^{2}$ has a middle term --- it is $7+2\sqrt{10}$, never $7$, for the same
reason $(x+4)^{2}\ne x^{2}+16$. The one fact that runs the whole day is
$\sqrt{a}\cdot\sqrt{a}=a$: \textbf{a square root times itself erases the radical}. That is what makes
$\left(4+\sqrt3\right)\left(4-\sqrt3\right)=16-3=13$ come out \emph{rational} --- the two middle
terms cancel and the radicals square away. Those two factors are called \textbf{conjugates}: same
terms, opposite middle sign (the same trick that turned $a+bi$ real in Lesson 2.4). Finally,
\textbf{rationalizing} clears a radical out of a denominator by multiplying top and bottom by the
same thing --- $\frac{\sqrt b}{\sqrt b}$ for a single radical, the \emph{conjugate} for a binomial.
It is always safe, because you are only multiplying by $1$.
\end{remindbox}

\end{document}
